\chapter{Conclusions and Future Work}
\label{cha:conclusioni}

\section{Conclusions}
\label{sec:conclusion}

This work presented a selective FPGA offload strategy for LFM2.5-1.2B, a model
co-designed for CPU inference~\cite{amini2025lfm2technicalreport}, targeting the MatMul +
SwiGLU feed-forward block to enable conversational audio inference pipelines on
resource-constrained edge platforms. The implementation integrates with the existing
llama.cpp/ggml flow and demonstrates measurable energy-efficiency gains, especially at low
and moderate thread counts, while preserving framework compatibility. Although the FPGA
path does not consistently exceed CPU Only throughput at higher thread counts, the results
clarify the limiting factors: practical on-device performance is jointly determined by
compute parallelism, weight/data movement, and host-device orchestration. Importantly, the
design already operates near the usable logic and routing budget of the target FPGA, so
further gains are less likely to come from naive parallel replication and more likely from
memory-path restructuring, overlap of transfers with compute, and tighter stage-level
balancing (notably in Stages~2a/2b and~5). Overall, the study confirms that selective
feed-forward acceleration is a viable and energy-effective step toward conversational-level
edge inference, and provides a concrete roadmap for the next optimizations in
LFM2.5-Audio-1.5B deployment.

\section{Future Work}
\label{sec:future}

The deployed design attains 61\% of the 3.6\,t/s design ceiling, with the gap attributed to
three quantified system-level constraints. Addressing these constraints defines the roadmap
for future work.

The design exhibits memory-bound behavior, achieving 38--41\% of peak DDR bandwidth at
34--35\% sustained compute utilization across the projection stages. This is consistent with
known bottlenecks in autoregressive decoding, where bandwidth utilization typically reaches
only 29--43\% even on higher-end platforms~\cite{FlightLLM}. Closing this performance gap
requires improved data movement strategies. Future work will therefore focus on
memory-centric optimizations, including tiling, double-buffering, and layout restructuring,
targeting higher sustained bandwidth utilization.

At the system level, host--accelerator overhead limits scalability. In the current design,
each generated token incurs a fixed synchronous handoff: the CPU programs the accelerator's
AXI-Lite control registers, issues a start signal, and blocks until a UIO interrupt is
delivered, a per-token overhead that does not amortize with sequence length and grows under
host-side scheduler pressure at higher thread counts. In contrast, other co-designed
frameworks achieve more stable scaling through tighter compiler--runtime
integration~\cite{cano}. Future work will investigate asynchronous execution and pipelined
scheduling to reduce synchronization overhead and improve utilization.

Extending hardware coverage beyond the MatMul + SwiGLU block to additional kernels is also a
key direction for improving end-to-end performance. This will require additional compute
fabric on the FPGA. Importantly, the design already effectively utilizes the available
resources on the KV260, and remaining inefficiencies stem primarily from memory and
system-level constraints, rather than insufficient optimization effort. Similar work pushing
the same platform to its limit required complex operator fusion and custom data arrangements
to maximize memory bandwidth~\cite{Li2025Pushing}.

Although the low DSP utilization (20\%) might suggest that mapping the 4-bit $\times$ INT8
inner product to DSP48E2 slices could free LUT fabric, this substitution was evaluated during
development and found to increase routing pressure: wider DSP fanout paths introduced
additional timing-critical nets that impeded timing closure at 250\,MHz. The DSP slice savings
were therefore insufficient to offset the routing overhead at the current operating point, and
the LUT-based multiply-accumulate path was retained.

Overall, future improvements will require co-optimization of data movement, scheduling, and
numerical representation, as well as evaluation on platforms with higher bandwidth and
resources, bridging the gap between lightweight selective acceleration and more general
co-design frameworks.
